
\section{1. Introduction}

\subsection{1.1 Motivation}

Full self-attention is $\mathcal{O}(L^2)$ in sequence length $L$, which 
becomes prohibitive for $L > 512$. Sliding-window attention reduces this to 
$\mathcal{O}(L \cdot W)$, but at a fundamental cost: tokens outside the window 
cannot attend to each other. Long-range dependencies must propagate through 
the residual stream over many layers.

\textbf{Hybrid attention} addresses this by combining a cheap local branch with 
a cheap global branch. The design question is: what should the global branch 
look at?

\begin{itemize}
\item \textbf{Sliding window + global tokens} (Longformer): the global branch attends
\end{itemize}
  to a fixed set of "global" tokens. Simple but limited.
\begin{itemize}
\item \textbf{Sliding window + compressed memory} (Infini-Attention): the global
\end{itemize}
  branch attends to a \textbf{learned matrix}. Powerful but adds $H \cdot d_c \cdot d_h$ 
  parameters.
\begin{itemize}
\item \textbf{Sliding window + sequence compression} (this work, OHA): the global
\end{itemize}
  branch attends to \textbf{$d_c$ summary slots} computed by soft-attention over 
  the sequence.

\textbf{OHA's key idea:} Rather than learning a fixed memory matrix, compute a 
\textbf{data-dependent compressed summary} of the current sequence. The compressed 
summary adapts to each input.

\subsection{1.2 Problem Statement}

Given a pretrained transformer $\mathcal{M}$ with $N_\ell$ attention layers, 
we seek a modified attention block $\mathcal{A}'$ such that:

1. $\mathcal{A}'$ reproduces the pretrained behavior at initialization.
2. $\mathcal{A}'$ reduces attention cost from $\mathcal{O}(L^2)$ to 
   $\mathcal{O}(L \cdot W + L \cdot d_c)$.
3. $\mathcal{A}'$ maintains access to global context via the compressed 
   summary.
4. $\mathcal{A}'$ works across four attention architectures.

We do \textbf{not} aim to reduce model size. OHA targets \textbf{inference efficiency} 
for long sequences.

\subsection{1.3 Our Approach: OHA}

OHA replaces MHSA with a two-branch block:

1. \textbf{Local branch:} sliding-window attention, $\mathcal{O}(L W)$.
2. \textbf{Global branch:} attention to $d_c$ compressed summary slots, 
   $\mathcal{O}(L d_c)$.
3. \textbf{Fusion:} per-head gate.

The compressed summary is computed by \textbf{soft-attention pooling}:

\[
C_K = \text{softmax}(W_k^\text{comp} X, \text{dim}=L) \cdot K \in \mathbb{R}^{B \times d_c \times d}
\tag{1.1}
\]

where $W_k^\text{comp} \in \mathbb{R}^{d \times d_c}$ produces pooling weights. 
Similarly for $C_V$.

The global branch attends to these compressed summaries:

\[
\text{GlobalAttn}_i = \text{softmax}\!\left(\frac{Q_i C_K^\top}{\sqrt{d_h}}\right) C_V
\tag{1.2}
\]

The two branches are fused via a per-head gate:

\[
\text{output} = \gamma \odot \text{LocalAttn} + (1 - \gamma) \odot \text{GlobalAttn}
\tag{1.3}
\]

with $\gamma = \sigma(W_g X) \in [0, 1]^{B \times H}$.

The gate bias is initialized to $2.0$, giving $\sigma(2) \approx 0.88$ -- 
favoring the local branch at start.

\subsection{1.4 Contributions}

\begin{itemize}
\item \textbf{A sequence-compressed hybrid attention block.} OHA combines local
\end{itemize}
  sliding-window attention with a global branch that attends to $d_c$ 
  data-dependent summary slots.
\begin{itemize}
\item \textbf{Soft-attention pooling.} The compression mechanism uses softmax
\end{itemize}
  attention over the sequence, producing input-adaptive summaries.
\begin{itemize}
\item \textbf{Per-head gating.} Each head learns its own local/global balance,
\end{itemize}
  initialized to favor local.
\begin{itemize}
\item \textbf{Multi-architecture support.} Same codebase handles flat, nested,
\end{itemize}
  causal, and GQA+RoPE models.
\begin{itemize}
\item \textbf{Honest reporting.} OHA increases parameter count; the "compression"
\end{itemize}
  refers to sequence length, not model size.

\subsection{1.5 Paper Organization}

Section 2 reviews related work. Section 3 provides background. Section 4 
presents OHA. Section 5 details multi-architecture integration. Section 6 
gives setup. Section 7 reports results. Section 8 analyzes contributions. 
Section 9 provides guidance. Section 10 lists limitations. Section 11 
concludes.



\subsection{Notation}
